\section*{Chapter \ref{ch:b1:electrophilic-additions} --- Alkenes and Alkynes: Electrophilic Additions}

\begin{solution}{exo:b1:electrophilic-additions:1}
2-Chloropropane; 2-bromo-2-methylpropane; 1-bromo-1-methylcyclohexane;
2-chlorobut-2-ene.
\end{solution}

\begin{solution}{exo:b1:electrophilic-additions:2}
1,2-Dibromoethane; \textit{trans}-1,2-dibromocyclohexane;
1,2-dibromopropane.
\end{solution}

\begin{solution}{exo:b1:electrophilic-additions:3}
Propan-2-ol; 2-methylpropan-2-ol; 1-methylcyclohexan-1-ol (Markovnikov in
each case).
\end{solution}

\begin{solution}{exo:b1:electrophilic-additions:4}
Propyne (\omterm{def:b1:electrophilic-additions:acetylide}{terminal alkyne}): \ce{CH3C#CH + NH2- -> CH3C#C- + NH3}. Ethanol:
\ce{CH3CH2OH + NH2- -> CH3CH2O- + NH3}. But-2-yne and propene have no
acidic enough hydrogen.
\end{solution}

\begin{solution}{exo:b1:electrophilic-additions:5}
The $\pi$ pair takes \ce{H+} on the \ce{CH2} carbon: tertiary cation
\ce{(CH3)3C+}; a water \omterm{def:b1:lewis-resonance-vsepr:lewis}{lone pair} bonds to it: \ce{(CH3)3COH2+}; a water
molecule takes a proton: \ce{(CH3)3COH}. In hot concentrated acid water is
scarce and the volatile alkene escapes: the reverse reaction, \omterm{def:b1:alcohol-activation:dehydration}{dehydration},
proceeds.
\end{solution}

\begin{solution}{exo:b1:electrophilic-additions:6}
Protonation of 2-methylpropene gives a tertiary cation, of propene a
secondary one. The rate-determining protonation has a \omterm{def:b1:elementary-steps:energy-profile}{transition state}
resembling the cation (Hammond): the more stable tertiary cation is
reached over a lower barrier.
\end{solution}

\begin{solution}{exo:b1:electrophilic-additions:7}
The bromonium ion bridges one face of the ring; bromide opens it from the
other face: the two bromines end \textit{trans}. The \textit{trans}
dibromide is \omterm{def:b1:stereochemistry-in-depth:stereocentre}{chiral} (no mirror plane), but bromide attacks the two carbons
equally: the two \omterm{def:b1:stereochemistry-in-depth:enantiomers}{enantiomers} form in equal amounts, and the product is
racemic, optically inactive.
\end{solution}

\begin{solution}{exo:b1:electrophilic-additions:8}
Protonation of C1 gives the secondary cation on C2, next to the tertiary
C3 bearing a hydrogen; that hydrogen moves with its pair to C2 (a
1,2-hydride shift), leaving a tertiary cation on C3, captured by
\ce{Cl-}: 2-chloro-2-methylbutane.
\end{solution}

\begin{solution}{exo:b1:electrophilic-additions:9}
Propyne + sodium amide: the propynide ion; with 1-bromopropane (SN2):
\ce{CH3C#C-CH2CH2CH3}, hex-2-yne. Propynide with benzaldehyde, then dilute
acid: \ce{C6H5CH(OH)C#CCH3}, 1-phenylbut-2-yn-1-ol.
\end{solution}

\begin{solution}{exo:b1:electrophilic-additions:10}
\cip{E}-but-2-ene: bromonium on one face, bromide from the other at C2
or at C3; both attacks give the same compound, with descriptors \cip{2R,3S}:
the \omterm{def:b1:stereochemistry-in-depth:enantiomers}{meso compound}, one stereoisomer. \cip{Z}-but-2-ene: attack at C2 gives
\cip{2S,3S} (or its mirror image from the other face), attack at C3
\cip{2R,3R}; equal amounts: two \omterm{def:b1:stereochemistry-in-depth:isomers}{stereoisomers}, a \omterm{def:b1:stereochemistry-in-depth:enantiomers}{racemic mixture}.
\end{solution}

\begin{solution}{exo:b1:electrophilic-additions:11}
In the unsymmetrical bromonium ion, the more substituted carbon bears more
of the positive charge (its C--Br bond is longer and weaker, like a nearly
tertiary or secondary cation). Water, in large excess, attacks that carbon
from the face opposite the bridge: 1-bromopropan-2-ol.
\end{solution}

\begin{solution}{exo:b1:electrophilic-additions:12}
$D = 1$: an alkene. 2-Methylbut-1-ene and 2-methylbut-2-ene both give the
tertiary cation, hence 2-bromo-2-methylbutane and 2-methylbutan-2-ol. Their
proton NMR spectra differ: two vinylic protons (\ce{=CH2}, singlets) for
the first, one (\ce{=CH-}, a quartet) for the second.
\end{solution}

\begin{solution}{pb:b1:electrophilic-additions:1}
\textbf{1.} On each carbon \ce{CH3} outranks H: \cip{Z} has the two methyls on
the same side, \cip{E} on opposite sides.
\textbf{2.} \omterm{def:b1:stereochemistry-in-depth:isomers}{Stereoisomers} that are not mirror images: \omterm{def:b1:stereochemistry-in-depth:enantiomers}{diastereomers}.
\textbf{3.} Rotation about C=C would break the $\pi$ bond, which costs far
more energy than collisions provide at room temperature.
\textbf{4.} \cip{E}: its methyl groups do not crowd each other.
\textbf{5.} \ce{C4H8 + Br2 -> C4H8Br2}.
\textbf{6.} Two equivalent stereocentres: at most four, in fact three
(\cip{2R,3R}, \cip{2S,3S} and the meso \cip{2R,3S}).
\textbf{7.} The $\pi$ pair attacks one Br atom of \ce{Br2}, the Br--Br pair
leaves as \ce{Br-}, and a \omterm{def:b1:lewis-resonance-vsepr:lewis}{lone pair} of the attacked bromine bonds to the
second carbon: a three-membered bromonium ion.
\textbf{8.} In the bridged ion every atom has an octet; an open secondary
\omterm{def:b1:electronic-effects:intermediates}{carbocation} would have a carbon with six electrons.
\textbf{9.} A \omterm{def:b1:lewis-resonance-vsepr:lewis}{lone pair} of \ce{Br-} bonds to one ring carbon; the C--Br bond
of the bridge breaks, its pair going to the bridging bromine.
\textbf{10.} The bridge occupies one face; as in SN2, the \omterm{def:b1:electronic-effects:nucleophile}{nucleophile} enters
opposite the bond that breaks.
\textbf{11.} An \omterm{def:b1:electrophilic-additions:halonium}{anti addition}.
\textbf{12.} Water would compete with bromide for the bromonium ion and give
a bromohydrin.
\textbf{13.} \cip{2R,3S} (or \cip{2S,3R}, the same compound).
\textbf{14.} The same compound.
\textbf{15.} In a \omterm{def:b1:stereochemistry-in-depth:isomers}{conformation} with the two bromines eclipsed, a mirror plane
passes between C2 and C3; the descriptors are opposite.
\textbf{16.} \cip{2R,3R} and \cip{2S,3S}.
\textbf{17.} Equal amounts: a \omterm{def:b1:stereochemistry-in-depth:enantiomers}{racemic mixture}.
\textbf{18.} Yes: the two diastereomeric alkenes give different products
(meso against racemic).
\textbf{19.} None: the two \omterm{def:b1:stereochemistry-in-depth:enantiomers}{enantiomers} cancel.
\textbf{20.} $5.60/56.0 = \qty{0.100}{mol}$ each.
\textbf{21.} $0.100 \times 0.85 \times 215.8 = \qty{18.3}{g}$ each.
\textbf{22.} \textbf{$[\alpha] = 0^\circ$ (\omterm{def:b1:stereochemistry-in-depth:enantiomers}{meso compound}), \qty{18.3}{g}}.
\end{solution}
